\documentclass[12pt]{article}
\usepackage[utf8]{inputenc}
\usepackage[T1]{fontenc}
\usepackage{amsmath,amssymb,amsthm}
\newcommand{\VV}{V}
\newcommand{\EE}{E}
\newcommand{\GG}{G}
\begin{document}
\title{Submission B solution to Problem 10}
\date{}
\maketitle

\section*{Original Problem}

Prove the following statement:

Let $\Gamma$ be a finite simplicial graph with $|\VV(\GG)|\geq 3$, and $(M_v,\varphi_v)$ be a separable tracial von Neumann algebra for each $v\in \VV$. If $\Gamma$ is irreducible (in the sense that there does not exists subgraphs $\GG_1,\GG_2$ such that $ \VV(\GG_1)\cup \VV(\GG_2)=\VV(\GG) $ and every $(v,u)\in \EE(\GG) $ for all $v\in\VV(\GG_1),u\in \VV(\GG_2)$), and each $M_v$ has a trace zero unitary, then the graph product von Neumann algebra $M$ is properly proximal.

\section*{Problem Solved}

This write-up proves the following \textbf{abstract DKEP finite-cover upgrade lemma}, not the full graph-product theorem stated above. Thus it is a relaxation/different result: it assumes the relevant DKEP boundary pieces, support projections, support sandwich inequality, and relative proper proximality as hypotheses, and proves only the final upgrade from relative proper proximality for finitely many pieces to absolute proper proximality.

Let \((M,\tau)\) be a separable tracial von Neumann algebra, let \(H=L^2(M,\tau)\), and let \(J\) be the modular conjugation, so that \(\rho(x)=Jx^*J\) denotes the right action of \(M\) on \(H\). Assume \(M\) is placed in the DKEP properly proximal framework: there is a unital von Neumann algebra \(\mathcal N\), called the normal-bidual compactification, containing the left copy of \(M\) and containing support projections for DKEP boundary pieces. A DKEP boundary piece \(X\subset B(H)\) means a hereditary \(C^*\)-subalgebra, invariant under the left and right \(M\)-actions, whose elements are treated as \(X\)-compact operators. Its DKEP support projection \(q_X\in\mathcal N\) is the projection obtained as the weak-* limit in \(\mathcal N\) of any contractive approximate identity of \(X\), equivalently the least projection supporting the weak-* closure of \(X\). The compact support \(q_K\in\mathcal N\) is the support projection of \(K(H)\).

Let \(n\ge 1\), let \(U_1,\dots,U_n\) be sets with finite-cover condition
\[
\bigcap_{i=1}^n U_i=\varnothing,
\]
and let
\[
\mathcal F=\left\{\bigcap_{i\in I}U_i:\varnothing\ne I\subseteq\{1,\dots,n\}\right\}\cup\{\varnothing\}.
\]
For every \(U\in\mathcal F\), suppose there is a DKEP boundary piece \(X_U\) with support projection \(q_U\in\mathcal N\). Assume \(q_{\varnothing}=q_K\). Assume also the support sandwich inequality
\[
q_Uq_Wq_U\le q_{U\cap W}
\]
in the positive order of \(\mathcal N\), for all \(U,W\in\mathcal F\).

For a boundary piece \(X\), define
\[
S_X(M)=\{T\in B(H):[T,\rho(x)]\in X\text{ for every }x\in M\}.
\]
A state \(\psi\) on \(S_X(M)\) is \(M\)-central if \(\psi(xT)=\psi(Tx)\) for all \(x\in M\), \(T\in S_X(M)\). The algebra \(M\) is properly proximal relative to \(X\) if there is no \(M\)-central state on \(S_X(M)\) whose restriction to the left copy of \(M\) is normal. The algebra \(M\) has no amenable direct summand if there is no nonzero central projection \(z\in Z(M)\) such that \(zM\) is amenable.

Use the following DKEP obstruction theorem as a black box: if \(M\) has no amenable direct summand and is not properly proximal, then there exists an \(M\)-central obstruction state \(\Omega\) on \(\mathcal N\), normal on \(M\), such that \(\Omega(q_K)=0\); moreover, if \(M\) is properly proximal relative to \(X_U\), then every such obstruction state satisfies \(\Omega(q_U)=1\).

Under these assumptions, if \(M\) has no amenable direct summand and is properly proximal relative to each \(X_{U_i}\), then \(M\) is properly proximal.

\section*{Proof}

All products and inequalities below are taken inside the normal-bidual compactification \(\mathcal N\). We deliberately do \textbf{not} assume that the projections \(q_U\) commute.

\subsection*{1. A state-support fact}

Let \(A\) be a unital \(C^*\)-algebra, let \(\phi\) be a state on \(A\), and let \(p\in A\) be a projection with \(\phi(p)=1\). Then
\[
\phi(a)=\phi(pap)
\qquad\text{for every }a\in A.
\]

Indeed, let \(e=1-p\). Then \(e\) is a projection and
\[
\phi(e)=\phi(1)-\phi(p)=0.
\]
By Cauchy–Schwarz for the positive functional \(\phi\),
\[
|\phi(ea)|^2\le \phi(e^*e)\phi(a^*a)=\phi(e)\phi(a^*a)=0,
\]
so \(\phi(ea)=0\) for all \(a\). Similarly,
\[
\phi(ae)=\overline{\phi(ea^*)}=0.
\]
Decomposing
\[
a=pap+pae+eap+eae,
\]
and applying \(\phi\), all terms except \(pap\) vanish. Hence \(\phi(a)=\phi(pap)\).

\subsection*{2. Two support-one projections force support one on the intersection}

Let \(\phi\) be a state on \(\mathcal N\). Suppose \(U,W\in\mathcal F\) and
\[
\phi(q_U)=1,\qquad \phi(q_W)=1.
\]
Set
\[
p=q_U,\qquad s=q_W,\qquad r=q_{U\cap W}.
\]
By the state-support fact applied to \(p\) and \(a=s\),
\[
\phi(s)=\phi(psp).
\]
Since \(\phi(s)=1\), we have
\[
\phi(q_Uq_Wq_U)=1.
\]
By the support sandwich inequality,
\[
q_Uq_Wq_U\le q_{U\cap W}=r.
\]
Because \(r\) is a projection, \(0\le r\le 1\), so
\[
1=\phi(q_Uq_Wq_U)\le \phi(r)\le 1.
\]
Therefore
\[
\phi(q_{U\cap W})=1.
\]

No commutativity was used: \(q_Uq_Wq_U\) is positive because
\[
q_Uq_Wq_U=(q_Wq_U)^*(q_Wq_U),
\]
and the argument uses only the compression \(q_Uq_Wq_U\), not the product \(q_Uq_W\).

\subsection*{3. Apply the obstruction theorem}

Assume, toward contradiction, that \(M\) is not properly proximal. Since \(M\) has no amenable direct summand, the DKEP obstruction theorem gives an \(M\)-central obstruction state \(\Omega\) on \(\mathcal N\), normal on \(M\), such that
\[
\Omega(q_K)=0.
\]

For each \(i=1,\dots,n\), \(M\) is properly proximal relative to \(X_{U_i}\). Hence the second part of the obstruction theorem gives
\[
\Omega(q_{U_i})=1
\qquad\text{for every }i=1,\dots,n.
\]

\subsection*{4. Iterate the two-support argument}

Define
\[
V_k=U_1\cap\cdots\cap U_k
\qquad (1\le k\le n).
\]
We prove by induction that
\[
\Omega(q_{V_k})=1
\qquad\text{for every }k.
\]

For \(k=1\), this is exactly \(\Omega(q_{U_1})=1\).

Assume \(\Omega(q_{V_k})=1\). Since \(\Omega(q_{U_{k+1}})=1\), the result of Section 2, applied with \(U=V_k\) and \(W=U_{k+1}\), gives
\[
\Omega(q_{V_k\cap U_{k+1}})=1.
\]
But
\[
V_k\cap U_{k+1}=V_{k+1}.
\]
Thus
\[
\Omega(q_{V_{k+1}})=1.
\]
By induction,
\[
\Omega(q_{V_n})=1.
\]

The finite-cover condition says
\[
V_n=U_1\cap\cdots\cap U_n=\varnothing.
\]
Therefore
\[
\Omega(q_{\varnothing})=1.
\]
By hypothesis,
\[
q_{\varnothing}=q_K.
\]
Hence
\[
\Omega(q_K)=1.
\]

This contradicts the obstruction-state condition \(\Omega(q_K)=0\). Therefore the assumption that \(M\) is not properly proximal is false. Consequently, \(M\) is properly proximal.

\end{document}
